\documentclass[10pt]{article}

% Specify the margins and text size.
\setlength{\textwidth}{6.4in}
\setlength{\textheight}{9.8in}
\setlength{\oddsidemargin}{0pt}
\setlength{\evensidemargin}{0pt}
\setlength{\topmargin}{0pt}
\setlength{\hoffset}{.05in}
\setlength{\voffset}{-1in}

\setlength{\parskip}{5pt}
\setlength{\parindent}{0pt}

% Load some fonts and symbol packages
\usepackage{latexsym}
\usepackage{pifont}       % contains 'star' symbol for counterinsurgency handbook title
\usepackage{yfonts} 
\usepackage{amsmath}
\usepackage{amsfonts}

\usepackage{graphicx}     % actually, this is loaded by pstricks
\usepackage[T1]{fontenc}
\usepackage{ifthen}
\usepackage{pstricks,pst-grad,pst-text,pst-node,multido,pst-plot,calc,pst-3dplot}
%\usepackage[all]{xy}
%\usepackage{animate}

% The hyperref package inserts links into the pdf.
\definecolor{MyLinkColor}{rgb}{.1,.2,1}
\definecolor{MyCiteColor}{rgb}{.1,1,.2}
\definecolor{MyURLColor}{rgb}{.4,.4,.4}
\usepackage[backref=true,pagebackref=false,hyperindex,colorlinks=true,
  linkcolor=MyLinkColor,urlcolor=MyURLColor]{hyperref}


% The tweaklist package is something I found on the web.  It provides a simple interface
% for making changes to spacing used in the itemize and enumerate environments.  Comment
% this out if you don't care to use tweaklist.
\usepackage{tweaklist}
\renewcommand{\itemhook}{\setlength{\parskip}{2pt}\setlength{\parsep}%
{1pt}\setlength{\topsep}{0pt}\setlength{\itemsep}{0pt}}

\newcommand{\U}{\underline{\hspace{5pt}}}

\usepackage{listings}
\newcommand{\Z}{\hphantom{0}}

\begin{document}
%\newcommand{\ticket}[2]{\psframe(0,0)(0.5,0.5)\psframe(0.5,0)(1,0.5)\rput(0.25,0.25){#1}\rput(0.75,0.25){#2}}
\pagestyle{empty}
\lstset{language=R, showspaces=false, showstringspaces=false}


\href{http://www.su.edu}{\includegraphics[height=1.75cm]{sulogo.eps}}
\vspace{-1.79cm}

{{\ }\hfill \small
\begin{tabular}{cl}
& Math 207\\
& Introduction to Statistics\\
\end{tabular}
}
\setlength{\baselineskip}{1.05\baselineskip}
\bigskip

\begin{center}
\textbf{\large  Independent Events}
\end{center}

Multiplication Rule:  
\begin{itemize}
\item $P(A\;\;\mbox{and}\;\;B)=P(A)\,P(B\,\vert\,A)$.
\item If $A$ and $B$ are independent, then $P(A\;\;\mbox{and}\;\;B)=P(A)\,P(B)$.
\end{itemize}
In particular, $A$ and $B$ are independent if and only if $P(A\vert B)=P(A)$.
This is true if and only if $P(B\vert A)=P(B)$.  In other words, knowing $A$ doesn't change
the chance of $B$; knowing $B$ doesn't change the chance of $A$.
\bigskip


\newcommand{\ticket}[2]{\psframe(0,0)(0.5,0.5)\psframe(0.5,0)(1,0.5)\rput(0.25,0.25){#1}\rput(0.75,0.25){#2}}

1. A six-sided die has the sides with 1, 2 and 3 painted blue.  The sides with 4, 5, and 6
are painted red.

\hspace{10pt} a) I roll the die and the top side shows blue.  What is the chance that the
 top side is a 2?
\vspace{.3in}

\hspace{10pt} b) I roll the die and the top side shows blue.  What is the chance that the
 top side is a 5?
 \vspace{.3in}

\hspace{10pt} c) I roll the die and the top side shows red.  What is the chance that the
 top side is a 2?
\vspace{.3in}

\hspace{10pt} d) Are the color shown and the number shown independent or dependent?
\vspace{.3in}

\hspace{10pt} e) The top side shows blue.
What is the chance that the top side shows an even number?
\vspace{.3in}


\hspace{10pt} f) The top side shows an even number.
What is the chance that the top side is blue?
\vspace{.3in}


\hspace{10pt} g) The top side shows an blue.
What is the chance that the top side shows an odd number?
\vspace{.3in}


\hspace{10pt} h) The top side shows an odd number.
What is the chance that the top side is blue?
\vspace{.3in}
\medskip

2. In the box shown below, each ticket should have two numbers:
\begin{center}
\begin{pspicture}(0,0)(15,0.5)
\rput(0,0){\ticket{1}{}}
\rput(2,0){\ticket{1}{2}}
\rput(4,0){\ticket{1}{2}}
\rput(6,0){\ticket{1}{3}}
\rput(8,0){\ticket{3}{1}}
\rput(10,0){\ticket{3}{2}}
\rput(12,0){\ticket{3}{}}
\rput(14,0){\ticket{3}{}}
\psline(-0.2,0.6)(-0.2,-0.2)(15.2,-0.2)(15.2,0.6)
\end{pspicture}
\end{center}

A ticket will be drawn at random.  Can you fill in the blanks so that 
the two numbers are independent?
\vspace{.5in}

3. In the box shown below, each ticket should have two numbers:
\begin{center}
\begin{pspicture}(0,0)(15,0.5)
\rput(0,0){\ticket{5}{}}
\rput(2,0){\ticket{5}{2}}
\rput(4,0){\ticket{5}{2}}
\rput(6,0){\ticket{5}{3}}
\rput(8,0){\ticket{3}{1}}
\rput(10,0){\ticket{3}{2}}
\rput(12,0){\ticket{3}{}}
\rput(14,0){\ticket{3}{}}
\psline(-0.2,0.6)(-0.2,-0.2)(15.2,-0.2)(15.2,0.6)
\end{pspicture}
\end{center}

A ticket will be drawn at random.  Can you fill in the blanks so that 
the two numbers are independent?
\vspace{1in}


\vfill
\eject
\end{document}

\href{http://www.su.edu}{\includegraphics[height=1.75cm]{sulogo.eps}}
\vspace{-1.79cm}

{{\ }\hfill \small
\begin{tabular}{cl}
& Math 207\\
& Introduction to Statistics\\
\end{tabular}
}
\setlength{\baselineskip}{1.05\baselineskip}
\bigskip

\begin{center}
\textbf{\large  Mutually-Exclusive Events}
\end{center}


Sum Rule:
\begin{itemize}
\item $P(A\;\;\mbox{or}\;\;B)=P(A)+P(B) - P(A\;\;\mbox{and}\;\;B)$
\item If $A$ ane $B$ are mutually exclusive, then $P(A\;\;\mbox{or}\;\;B)=P(A)+P(B)$
\end{itemize}


\vfill
\eject

\end{document}

\#1 p.\ 213. Find the regression equation for predicting the final exam score from
the midterm score based on the following information:
\begin{center}
\begin{tabular}{rcc}
average midterm score $=$ 70, & SD $=$ 10\\
average final score $=$ 55, & SD $=$ 20, & $r=0.60$
\end{tabular}
\end{center}

